\section{Conclusion}

This paper investigated a central deployment problem in vehicular edge computing: whether deadline-compliant task delivery can be preserved as increasing numbers of vehicles simultaneously compete for heterogeneous communication links and shared MEC resources. MA-ROPE addresses this problem by jointly considering RAT selection, partial computation offloading, shared RSU queue contention, and federated cooperative learning across LTE, ITS-G5, WiFi-6, and 5G-NR.

The experiments demonstrate that successful packet reception alone is insufficient to characterize V2X offloading performance under density. At $N=32$ under URLLC stress, MA-ROPE achieves a fleet OTR of 0.834 compared with 0.525 for single-RAT AllOffload, while reducing energy consumption from 0.944 J to 0.706 J, a 25.2\% reduction. The distinction becomes particularly important for safety traffic: MA-ROPE achieves safety $\text{OTR}=0.752$, whereas AllOffload falls to 0.287, despite both methods retaining comparatively high packet reception ratios.

To characterize this scaling behavior directly, we introduced the Deadline Retention Ratio $R_D$. Across a fourfold increase from $N=8$ to $N=32$, MA-ROPE retains $R_D=0.940$ of its reference-density on-time delivery performance, whereas single-RAT AllOffload retains $R_D=0.599$. This result captures the principal benefit of heterogeneous queue-aware offloading: MA-ROPE does not merely achieve useful deadline performance at one operating point; it preserves substantially more of that performance as contention increases.

The results also expose the trade-off between centralized performance and deployability. MADDPG achieves a $U_{OTR}$ advantage of approximately 0.009 over MA-ROPE at $N=32$ and lower mean latency, but requires centralized observation sharing and consumes approximately 34.0\% more energy. MA-ROPE therefore provides a different operating point: near-upper-bound deadline utility with lower energy expenditure and without requiring centralized collection of raw vehicle observations.

The additional decision complexity does not translate into a significant runtime penalty. The measured simulator-step time is 166.6 \textmu s for the multi-RAT configuration versus 167.6 \textmu s for single-RAT offloading, while the complete N=8 inference-plus-environment step requires approximately 644.3 \textmu s. These measurements remain well within the evaluated control-time budget and indicate that heterogeneous decision making itself is not the computational bottleneck.

Several limitations remain. The deep safety validation at $N=8$ is MA-ROPE-specific and should not be directly compared with the multi-method $N=32$ safety experiment. The current evaluation also uses synthetic mobility and task distributions, while FedAvg reduces raw-observation sharing but does not itself provide formal privacy guarantees. Future work should therefore investigate non-IID vehicle and task distributions, trace-driven mobility, RSU handover and failure, secure aggregation or differential privacy, and sensitivity to utility weights and MEC service rates.

Overall, the results support two conclusions. First, heterogeneous RAT selection, partial offloading, and shared MEC contention should be optimized jointly when V2X workloads operate under strict deadlines. Second, evaluation at a single vehicle density can conceal the principal failure mode of aggressive offloading. Deadline-aware fleet scaling---and the proposed $R_D$ measure of deadline retention---provides a more direct indication of whether an offloading policy remains useful as the deployment becomes dense.
